\adnote{This chapter description from the proposal will function as the "roadmap" for the structure and contents of this dissertation chapter.}

%========================
% Chapter 2 — Literature Review (Proposal)
%========================

% \chapter{Literature Review} % uncomment if your class uses \chapter

This chapter establishes the theoretical foundations for \textit{Speculocultural Technopoiesis} (ST) by synthesizing work in philosophy of technology, Black studies and speculative traditions, speculative/critical design, performance and sonic studies, computational culture, and research-creation. It advances a reading of mediation as culturally specific, identifies where Eurocentric defaults are embedded in audio tools, and articulates how Black sonic epistemologies can be made computationally operational.

\subsubsection{Objectives}
\begin{itemize}[noitemsep]
	\item Establish mediation as culturally specific and identify where Eurocentric defaults reside in audio tools.
	\item Articulate Black speculative and sonic traditions as sources of computational logic (not metaphor).
	\item Frame algorithms and interfaces as rhetorical/cultural, motivating abolitionist redesign.
	\item Ground a hybrid method (critical technical practice $+$ research-creation $+$ autoethnography/co-design) that operationalizes ST.
	\item Define evaluation criteria centered on embodied cultural authenticity and transmissible documentation.
\end{itemize}

\subsubsection{Intended Sections/Subsections}
\begin{itemize}[nosep]
	%-----------------------------------------------------------------------
	\item \subsubsection{Technologies as Cultural Mediators (Unsettling Foundations)}
	\begin{itemize}[nosep]
		\item \subsubsection{Post-Phenomenology (Ihde)}
			\textbf{Core claim.} Technologies are non-neutral mediators of perception/action; human–technology relations (embodiment, hermeneutic, alterity, background) co-constitute experience. \\
			\textbf{Tension.} Emphasis on embodiment can under-specify cultural specificity of mediation. \\
			\textbf{Implication for ST.} Treat DAW defaults, time bases, and representational schemas as culturally loaded mediations to be surfaced and re-authored.
		\item \subsubsection{Agential Realism (Barad)}
			\textbf{Core claim.} Phenomena emerge through \textit{intra-action} within apparatuses; apparatus design is material–discursive. \\
			\textbf{Tension.} Needs translation into concrete interface/algorithm design practices. \\
			\textbf{Implication for ST.} Consider audio interfaces as apparatuses that can be refactored to diffract, not reflect, Eurocentric norms.
		\item \subsubsection{Technopoiesis (Guerra \& Ostergaard)}
			\textbf{Core claim.} Technology emerges via feedback across exigencies, representations, materials, and cultural entrenchment. \\
			\textbf{Implication for ST.} Use Black sonic epistemologies as inputs guiding representational forms and material implementations (algorithms).
		\item \subsubsection{Platform/Material Studies (Montfort \& Bogost; Terranova; Malazita)}
			\textbf{Core claim.} Platforms and media ecologies encode assumptions; infrastructures shape possible practices. \\
			\textbf{Implication for ST.} Read audio stacks (plugins, engines, controllers) as sites where defaults can be strategically re-authored.
		\item \subsubsection{Synthesis} Mediation is culturally specific; the audio apparatus/platform encodes it. ST targets those encodings.
	\end{itemize}
	%---------------------------------------------------------------------
	\item \subsubsection{Speculative Worlds and Futures (Remixing Epistemologies)}
	\begin{itemize}[nosep]
		\item \subsubsection{Critical \& Speculative Design (Dunne \& Raby)}
		\textbf{Use.} Design as method to open alternative worlds and question defaults. \\
		\textbf{Gap.} Risks aesthetics of speculation without material reconfiguration.
		\item \subsubsection{Black Speculative Traditions (Nelson; Womack; Anderson \& Jones; Eshun)}
		\textbf{Use.} Afrofuturism and sonic fiction articulate survival, liberation, and worlding through sound. \\
		\textbf{Bridge to computation.} Treat grooves, polyrhythm, and call–response as design logics (not metaphors).
		\item \subsubsection{Critical Fabulation (Hartman)}
		\textbf{Use.} Speculative method to address archival silences ethically. \\
		\textbf{Bridge.} Guides documentation and naming practices for culturally grounded tools.
		\item \subsubsection{Synthesis} Speculation becomes design requirement when tied to sonic logics and tool modification.
	\end{itemize}
	%---------------------------------------------------------------------
	\item \subsubsection{Fugitivity and Resistant Practice}
	\begin{itemize}[nosep]
		\item \subsubsection{Moten \& Harney — The Undercommons}
		\textbf{Use.} Fugitivity and study as collective practices of refusal and invention. \\
		\textbf{Bridge.} “Unauthorized uses” as deliberate interface repurposing.
		\item \subsubsection{Wynter — Sociogenic Critique}
		\textbf{Use.} Unsettling the coloniality of “Man”; plural epistemologies of the human. \\
		\textbf{Bridge.} Grounds “specific-to-universal” stance: begin in specificity to revise defaults.
		\item \subsubsection{Robinson; Kelley — Radical Imagination}
		\textbf{Use.} Genealogies of Black radical thought; freedom dreaming. \\
		\textbf{Bridge.} Frames design aims as abolitionist re-worlding in code and interface.
		\item \subsubsection{Synthesis} Resistance is method: encode refusals as computational affordances.
	\end{itemize}
	%---------------------------------------------------------------------
	\item \subsubsection{Embodiment and Memory (Embodying Knowledge)}
	\begin{itemize}[nosep]
		\item \subsubsection{Performance as Epistemology (Johnson; DeFrantz; Black Performance Theory)}
		\textbf{Use.} Performance produces knowledge; memory is enacted somatically. \\
		\textbf{Bridge.} Validate design via embodied cultural fit, not only task efficiency.
		\item \subsubsection{Sonic Traditions (Gilroy; Weheliye)}
		\textbf{Use.} Sound as diasporic technology; assemblage, groove, and viscus. \\
		\textbf{Bridge.} Derive time bases, interaction grammars, and data structures from sonic practice.
		\item \subsubsection{Embodied Interaction \& Enchantment (Höök; Jones; Sha Xin Wei)}
		\textbf{Use.} Somaesthetic and multisensory interaction; poiesis in material practices. \\
		\textbf{Bridge.} Privilege movement/sound-first interaction over ocularcentric control.
		\item \subsubsection{Sonic Interaction Design (Gaver; Rocchesso; Serafin; Franinović)}
		\textbf{Use.} Sound as primary design material for interaction. \\
		\textbf{Bridge.} Provide techniques to prototype/assess sonic affordances in practice.
		\item \subsubsection{Synthesis} Embodied validation and sonic logics anchor how tool changes are judged.
	\end{itemize}
	%---------------------------------------------------------------------
	\item \subsubsection{Computational Culture and Algorithmic Power}
	\begin{itemize}[nosep]
		\item \subsubsection{Critical Code \& Software Studies (Marino; Fuller)}
		\textbf{Use.} Code/software as cultural text; read defaults critically. \\
		\textbf{Bridge.} Expose which abstractions (e.g., 4/4 grids) are not neutral.
		\item \subsubsection{Procedural Rhetoric (Bogost)}
		\textbf{Use.} Processes make arguments; algorithms persuade by what they enable/constrain. \\
		\textbf{Bridge.} Treat algorithmic time and interaction grammars as rhetorical.
		\item \subsubsection{Discriminatory Design \& Abolitionist Tools (Benjamin)}
		\textbf{Use.} New Jim Code; discriminatory defaults in everyday tech. \\
		\textbf{Bridge.} Pursue “abolitionist tools” by encoding Black sonic values.
		\item \subsubsection{Synthesis} Defaults are arguments; ST supplies counter-arguments as computational designs.
	\end{itemize}
	%---------------------------------------------------------------------
	\item \subsubsection{Methodological Precedents (Toward Research-Creation)}
	\begin{itemize}[nosep]
		\item \subsubsection{Critical Technical Practice (Agre; Sengers)}
		\textbf{Use.} Keep critique and making in continuous loop; reflect on hidden assumptions during design. \\
		\textbf{Bridge.} Operational template for ST’s cycle.
		\item \subsubsection{Critical Making (Ratto; Rosner; Hertz)}
		\textbf{Use.} Material critique that produces knowledge through artifacts. \\
		\textbf{Bridge.} Aligns with “Remixing” as embodied modification.
		\item \subsubsection{Autoethnography \& Collaborative Autoethnography (Ellis; Chang; Denzin)}
		\textbf{Use.} Legitimate practice-based, situated knowledge; enable cohort reflection. \\
		\textbf{Bridge.} Grounds \textit{Intimate Witnessing} and documentation.
		\item \subsubsection{Participatory/Feminist HCI \& Workshops (Schuler \& Namioka; Bardzell; Andersen \& Wakkary)}
		\textbf{Use.} Situated, reflexive co-design; magic-machine workshops as generative structures. \\
		\textbf{Bridge.} Provides session format for small-cohort sonic design.
		\item \subsubsection{Research-Creation (Loveless; Manning \& Massumi; Bolt)}
		\textbf{Use.} Making as theory; concepts emerge through practice. \\
		\textbf{Bridge.} Validates case studies as sites of conceptual production.
	\end{itemize}
	%-------------------------------------------------------------------
	\item \subsubsection{Synthesis: Toward Speculocultural Technopoiesis}
	\textbf{Integrated stance.} (1) Mediation is culturally specific and materially encoded; (2) Black speculative and performance traditions provide computational logics; (3) Defaults are rhetorical and often discriminatory; (4) Research-creation and critical technical practice supply the means to modify apparatuses ethically and collaboratively. \\
	\textbf{Consequence.} ST formalizes a cycle (Unsettling $\rightarrow$ Remixing $\rightarrow$ Embodying $\rightarrow$ Transmitting) and four territories (Fugitive Speculation, Material Memory, Sensory Worlding, Intimate Witnessing) as a reproducible method for culturally grounded tool modification and redesign.	
\end{itemize}

% --- Chapter 2 Key References (keyword-based filter) ---
\begin{comment}
\begin{refsection}[references.bib] % ensure your ch.2 items have keyword=ch2
	\nocite{*}
	\printbibliography[heading=subbibliography,title={Key References}]
\end{refsection}
\end{comment}